\documentclass[11pt]{article}
\usepackage[margin=1in]{geometry}
\usepackage{amsmath,amssymb}

\title{PDE-Oriented Notes for CV-DV LCHS Heat-Equation Runs}
\date{\today}

\begin{document}
\maketitle

\section*{1. Target Equation}
We target the original discrete heat-equation map
\begin{equation}
u_{\mathrm{PDE}}(t)=e^{-\alpha t T}u(0),
\end{equation}
with fixed $T$ from the 2-qubit finite-difference decomposition:
\begin{equation}
T = 2(I\otimes I) - (I\otimes X) - \frac{1}{2}(X\otimes X + Y\otimes Y).
\end{equation}
In the current project mode, no extra operator scaling/shift is used.

\section*{2. CV Preparation Parameters That Matter}
The non-Gaussian LCHS coefficients are determined by
\begin{equation}
C_n \propto \int dk\;
e^{-\gamma k^2}\,g_{\beta_k}(k)\,
\frac{H_n(k/e^{r'})}{\sqrt{2^n n!}},
\qquad
\gamma=e^{-2r'}-e^{-2r},\quad r'<r,
\end{equation}
with kernel
\begin{equation}
g_{\beta_k}(k)=\frac{\exp((1+ik)^{\beta_k})}{1-ik}.
\end{equation}
Hence the relevant sweep knobs are:
\begin{itemize}
\item $r=r_{\texttt{target}}$ (post-selection squeezing),
\item $r'=r_{\texttt{prime}}$ (prep-basis squeezing),
\item $\beta_k=\texttt{kernel\_beta}$ (kernel shape).
\end{itemize}

\section*{3. Metrics for Correctness}
Primary metric:
\begin{equation}
\varepsilon_{\mathrm{PDE}}
=\frac{\|u_{\mathrm{PDE}}-\sqrt{p_{\mathrm{post}}}\,\psi_{\mathrm{principal}}\|}{\|u_{\mathrm{PDE}}\|},
\end{equation}
where $\psi_{\mathrm{principal}}$ is the principal eigenvector of the reconstructed post-selected DV density matrix.

Secondary metric: post-selection probability $p_{\mathrm{post}}$.

Diagnostic metrics: mixed-state fidelity
\begin{equation}
F_{\mathrm{mix}}=\langle \hat u_{\mathrm{PDE}}|\rho_{\mathrm{post}}|\hat u_{\mathrm{PDE}}\rangle,
\end{equation}
and purity $\mathrm{Tr}(\rho_{\mathrm{post}}^2)$.

\section*{4. Theory-Guided Optimization Model}
We optimize over
\begin{equation}
\theta=(r,r',\beta_k),
\end{equation}
with physically consistent inequality
\begin{equation}
r' < r.
\end{equation}

Useful constrained form:
\begin{equation}
\min_{\theta}\ \varepsilon_{\mathrm{PDE}}(\theta)
\quad\text{s.t.}\quad
p_{\mathrm{post}}(\theta)\ge p_{\min},\ 
\gamma(\theta)\ge \gamma_{\min},\ 
n_{\mathrm{eff}}(\theta)\le n_{\mathrm{eff,max}}.
\end{equation}

Here
\begin{equation}
\gamma(\theta)=e^{-2r'}-e^{-2r},
\qquad
n_{\mathrm{eff}}(\theta)=\frac{1}{\sum_n |C_n(\theta)|^4}.
\end{equation}

The implemented unconstrained penalty objective is
\begin{equation}
\begin{aligned}
J(\theta)=\varepsilon_{\mathrm{PDE}}(\theta)
&+\lambda_p\frac{[p_{\min}-p_{\mathrm{post}}(\theta)]_+}{p_{\min}}
+\lambda_{\gamma}\frac{[\gamma_{\min}-\gamma(\theta)]_+}{\gamma_{\min}} \\
&+\lambda_n\frac{[n_{\mathrm{eff}}(\theta)-n_{\mathrm{eff,max}}]_+}{n_{\mathrm{eff,max}}}
-10^{-3}p_{\mathrm{post}}(\theta),
\end{aligned}
\end{equation}
with $[x]_+ = \max(0,x)$.

\section*{5. Practical Rule}
When selecting final settings, rank by:
\begin{enumerate}
\item smallest $\varepsilon_{\mathrm{PDE}}$,
\item largest $p_{\mathrm{post}}$ as tie-break.
\end{enumerate}
Use fidelity only as a consistency diagnostic, not as the primary objective.

\end{document}
